\section{Related Work}

\paragraph{Rule-governed, evidence-bearing document review.}
Document inspection has a long history of using explicit review procedures and checklists~\cite{fagan1976,lait2000}. Requirements engineering adds quality criteria, smell detection, and traceability to this tradition~\cite{iso29148,genova2013,femmer2017}. We use \emph{rule-governed, evidence-bearing document review} to denote a task in which an organisation supplies review requirements, a system associates proposed findings with evidence in the submitted material, and an accountable reviewer resolves issues that the automated procedure cannot settle. Requirements specifications, policy-governed contracts, and technical solution proposals provide different instances of this task~\cite{chinnam2025,alsawalqah2026,malik2026}. Security requirements and industrial certification research further illustrate the importance of explicit review criteria~\cite{secreq2025,cyberreq2026}, while traceability has also been studied for an in-the-wild augmented-reality learning service~\cite{trace2026}. These sources motivate the broader task setting; their existence does not establish that HARK transfers across those domains. Technical solution proposal review remains the only empirical case in this study.

\paragraph{Requirements quality and operational verifiability.}
Requirements-smell analysis detects cues such as weak wording, placeholders, and missing units~\cite{femmer2017}. Veizaga et al. combine smell detection with recommendations based on pre-trained language models~\cite{smell2024}, and AirReq applies smell detection and elimination to commercial aircraft system requirements~\cite{airreq2025}. HARK draws on these techniques for structural checks and candidate localisation, but a surface cue is not by itself a complete assessment of whether a requirement is satisfied. A further distinction concerns the meaning of \emph{verifiability}. As a requirements-quality attribute, verifiability concerns whether satisfaction can be established through an appropriate verification procedure~\cite{iso29148,genova2013}; it does not mean that the current document parser can decide the requirement. HARK's tiers are operational assignments relative to the available parser, verifier, representation, and rule package. A requirement requiring measurements, external records, or an image may be verifiable in principle while remaining outside this prototype's deterministic checking capability. The tier therefore specifies output eligibility within an implementation, not an intrinsic quality ranking of the requirement.

\paragraph{LLM-assisted requirements review and evaluation.}
Reviews of LLM applications in requirements engineering describe uses across extraction, analysis, and specification~\cite{hemmat2025}; a dual-perspective review distinguishes using LLMs for code verification from verifying their generated outputs~\cite{dolcetti2025}. MARV uses multiple role agents and retrieval of standards for software requirements specification validation~\cite{chinnam2025}. ReqPlan-Eval explicitly couples evidence-aware routing with human requirements-to-planning decisions and discloses the limits of its evaluation~\cite{alsawalqah2026}. These approaches supply relevant precedents for evidence use and reviewer involvement. The inspected MARV material does not establish the detailed behaviour of document-range validation or persistent feedback, so we do not infer that those capabilities are absent.

Hallucination research motivates checking generated claims~\cite{huang2025hall}, but its findings should not be converted into a universal claim that models cannot perform deterministic tasks. Long-context position effects~\cite{liu2024lost}, limits of intrinsic reasoning self-correction without external feedback~\cite{huang2024self}, and model-contingent polarity in a hotel-review annotation task~\cite{polarity2026} concern distinct failure modes and settings. They motivate separating model judgement from permission to issue a result, rather than proving the effectiveness of HARK's particular routing policy. Likewise, the evaluation literature~\cite{chang2024eval} motivates reporting the corpus, judge, and task conditions under which a mechanism was tested. HARK's deterministic-judge evaluation addresses this bounded mechanism question, not the capabilities of arbitrary deployed LLMs.

\paragraph{Neuro-symbolic policy review.}
Neuro-symbolic integration offers several ways for symbolic components to constrain learned components~\cite{garcez2023}. PolicyGuard compiles organisational policies into typed relational rules, uses LLM extraction to ground atoms from evidence, and applies symbolic rule evaluation~\cite{malik2026}. Its v1 report also documents legal-expert refinement during review-engine construction and validation: Appendix G describes editing extraction questions, correcting rule-combination logic, and treating a triggered finding as a case-specific exception while preserving the audit trail. It is therefore inaccurate to characterise PolicyGuard as having no feedback or human maintenance. The comparison instead concerns the role of the evidence-producing component and the lifecycle of review knowledge: PolicyGuard reports atom grounding, symbolic evaluation, and local expert corrections; HARK studies rule-specific eligibility for automatic findings and a stored review-feedback layer. These are different architectural choices, not evidence of a general deficit in PolicyGuard or proof that HARK has implemented complete human-controlled evolution.

Galitsky and Rybalov study symbolic checking of LLM-generated fault explanations and control actions against physical and sensor constraints in \emph{process control}~\cite{neuro2026}. That setting illustrates constraint-based rejection, but it is not a study of design-document review. Nor do these examples justify classifying all neuro-symbolic systems as assigning the same authority to every learned output. HARK's relevant design choice is to make the eligible output and its evidence conditions explicit for each rule before reviewing a document.

\paragraph{Industrial review and editable feedback knowledge.}
Code review provides adjacent evidence about learning from reviewer reactions, although its artefacts and evaluation units differ from document review. Pornprasit and Tantithamthavorn compare fine-tuning and prompting for code review automation~\cite{pornprasit2024}. BitsAI-CR's v1 architecture contains \emph{two fine-tuned LLMs}, RuleChecker and ReviewFilter; RuleChecker is not a deterministic rule engine. Its data flywheel updates review rules and their taxonomy as well as supplying data for model improvement~\cite{sun2025}. The RovoDev report evaluates a deployed reviewer using offline and online measurements and developer feedback~\cite{rovo2026}. Lin et al. analyse developer replies to CodeRabbit reviews, reporting rejection and misalignment within the reply-bearing sample rather than establishing rejection rates for all generated comments~\cite{lin2026}.

Feedback is not necessarily hidden in model weights or inaccessible to users. CodeRabbit's official Learnings documentation, as accessed on 5 November 2026, describes stored user-generated learnings that can be viewed and edited, with role-governed deletion through the interface and removal requests through chat~\cite{coderabbitLearnings}. These controls are a positive precedent for inspectable review knowledge. They do not, by themselves, establish complete historical rollback, permission-ceiling enforcement, or replay of the exact knowledge content used for an earlier review. Conversely, those capabilities cannot be declared absent merely because an industrial evaluation paper did not assess them. HARK's stored feedback should be compared with these concrete mechanisms, not with a homogeneous category of supposedly opaque industrial agents.

\paragraph{Grounding, non-detection, and pending decisions.}
Attributed generation links claims to cited sources~\cite{gao2023}. Claim-level evaluation of verbatim citation distinguishes quotation compliance from whether the quoted passage actually supports the claim~\cite{zhang2026cite}. HARK adds a rule-associated, tool-side eligibility check: the evaluated prototype locates a submitted quotation and checks overlap with verifier-provided candidate evidence for that rule. This \emph{located-overlap gate} is narrower than a guarantee of complete span containment or semantic entailment. A located quotation can still be irrelevant or misinterpreted, and candidate retrieval can miss relevant material.

Non-detection also needs an explicit evidence boundary. A verifier-owned non-detection record reports that a specified procedure did not find expected material in its declared checking scope; it is not proof that the requirement is unsatisfied in the world, or that relevant content is absent from every representation of the document. HARK reserves such records for verifier output rather than accepting a model-generated assertion of absence. When a hybrid rule lacks usable candidates, or a proposed finding does not meet its evidence conditions, an explicit \emph{pending} state preserves the unresolved obligation instead of turning it into either a pass or a confirmed failure. Documents remain untrusted inputs: prompt-injection attacks against LLM judges~\cite{shi2024inject} motivate this separation, but an evidence gate alone does not establish resistance to injection or protect every write path.

\paragraph{Human control, memory, and version requirements.}
Human-in-the-loop research distinguishes learning arrangements according to how control is allocated between system and expert~\cite{mosqueira2023}. HARK adopts this control question at the level of review rules and feedback events; the presence of a named reviewer or a stored approval text is not equivalent to authenticating the approving human. Agent-memory research~\cite{wang2024agents}, self-refinement~\cite{madaan2023}, reflection-based memory~\cite{shinn2023}, and self-modifying agents~\cite{zhang2025dgm} show multiple ways to reuse interactions or improve behaviour. Neither prompting nor retrieval inherently precludes audit logs, persistent consolidation, or versioned knowledge. Those properties depend on the surrounding storage, update, and permission mechanisms.

Memory extraction can expose sensitive stored content~\cite{mempriv2025}. HARK's structured feedback records and short exemplars make stored fields explicit, but this representation has not been shown to provide a privacy guarantee. More generally, bounded evolution requires separately establishing the identity and authority of the approving actor, rejecting updates that exceed an authorised permission ceiling, and retaining the actual content needed to replay a prior knowledge state. The evaluated prototype provides operational tier routing, a located-overlap gate, and stored feedback; \emph{complete human authentication, permission-ceiling enforcement, and content replay are not yet implemented}. Consequently, feedback storage and version labels are reported mechanisms, not evidence that every update is human-authorised, reversible, or beneficial over time.

\paragraph{Positioning and contribution boundary.}
Table~\ref{tab:positioning} compares concrete reports or explicitly labelled family illustrations, rather than scoring entire research traditions. The positive contribution studied here is a system representation that connects \emph{rule-specific output eligibility}, \emph{verifier-owned non-detection records}, \emph{explicit pending outcomes}, and \emph{feedback/version requirements} within one review workflow. The three tier labels are not themselves a novelty claim. We neither claim a first combination based on unverified negative capabilities of other systems nor claim that the full human-gated lifecycle has been achieved. The measurements reported later concern the implemented proposal-review prototype and its stated judge; broader document classes and strict human-authorised evolution remain outside the demonstrated scope.

\begin{table}[!ht]
\caption{Mechanism-level positioning using named reports and versions. ``Not established'' denotes a limit of the inspected evidence, not an assertion that a system lacks the capability. Rows are not a common-benchmark comparison; reported controls are not treated as proof of complete rollback or human authentication.}
\label{tab:positioning}
\centering
\footnotesize
\setlength{\tabcolsep}{3pt}
\begin{tabularx}{\textwidth}{@{}>{\raggedright\arraybackslash}p{2.6cm}>{\raggedright\arraybackslash}X>{\raggedright\arraybackslash}X>{\raggedright\arraybackslash}X@{}}
\toprule
Report or illustration & Knowledge and decision mechanism & Evidence and unresolved outcomes & Feedback and demonstrated control boundary \\
\midrule
2025 keyword prototype (reimplemented baseline) & Substring checks per proposal rule & Keyword matches; no satisfaction semantics & No persistent feedback in this baseline implementation \\
Requirements-smell family illustration~\cite{femmer2017,smell2024,airreq2025} & Cue detection; some systems add learned recommendations & Requirement-level quality cues; coverage depends on the analyser & No uniform family-wide lifecycle or rollback claim \\
MARV, 2025 thesis~\cite{chinnam2025} & Role agents with retrieved standards & Requirements validation; detailed document-range gate not established & Persistent feedback and rollback not established by inspected material \\
ReqPlan-Eval, 2026 report~\cite{alsawalqah2026} & Evidence-aware routing; final human decisions & Evidence records and disclosed evaluation limits & Knowledge evolution not reported as a tested mechanism \\
PolicyGuard, arXiv v1~\cite{malik2026} & Typed relational policy; LLM-grounded atoms; symbolic evaluation & Atom-level evidence and rule traces & Construction/validation expert corrections and audit trail; complete deployed rollback not established \\
BitsAI-CR, arXiv v1~\cite{sun2025} & Fine-tuned LLM RuleChecker and ReviewFilter & Generated comments and learned secondary validation & Data flywheel refines models and rule taxonomy; per-reviewer rollback not established \\
RovoDev, arXiv v2 report~\cite{rovo2026} & Deployed code-review generation and filtering & Comment quality evaluated offline and online & Developer feedback evaluated; full knowledge-history replay not established \\
CodeRabbit Learnings, official documentation accessed 5 Nov.\ 2026~\cite{coderabbitLearnings} & Stored user-generated review knowledge & Knowledge applied to subsequent code reviews & View/edit/delete controls; these alone do not establish complete historical rollback \\
\textbf{HARK evaluated prototype} & Operational tier routing and rule-specific output eligibility & Located-overlap gate; verifier non-detection; explicit pending & Stored feedback; complete human authentication, permission ceiling, and content replay not yet implemented \\
\bottomrule
\end{tabularx}
\end{table}


